% SLICECANCELDEPTH1 2026-10-18 02:05 Asia/Seoul
\documentclass[11pt]{article}
\usepackage[T1]{fontenc}
\usepackage{lmodern,amsmath,amssymb,amsthm,booktabs,hyperref}
\usepackage[margin=0.9in]{geometry}
\newtheorem{theorem}{Theorem}
\newtheorem{lemma}[theorem]{Lemma}
\newtheorem{corollary}[theorem]{Corollary}
\newcommand{\EXACT}{\textbf{PROVED/EXACT}}
\newcommand{\KNOWN}{\textbf{KNOWN THEOREM}}
\newcommand{\COMP}{\textbf{COMPUTATIONAL}}
\newcommand{\OPEN}{\textbf{OPEN}}
\newcommand{\RET}{\textbf{FALSIFIED/RETIRED}}
\title{SLICECANCELDEPTH1:\\Respect-Closure Barriers and Exponential Inclusion--Exclusion Cancellation Depth}
\author{P versus NP project}
\date{18 October 2026}
\begin{document}\maketitle

\noindent\textbf{Status.} The respect-closure theorem, the cardinality-collapse corollary for the full exact-density UNSAT positive family, the factorial-moment recurrence, and the exponential Bonferroni cancellation-depth theorem are \EXACT at their stated models, pending independent review. The second factorial-cumulant exponent separation is \EXACT at the exponential-rate lower-bound level. Higher fixed-order cumulant exponents are \COMP. No unrestricted set-circuit lower bound is proved; $P\ne NP$ remains \OPEN.

\section{IDEAS PROPOSED BEFORE LITERATURE}
\paragraph{A: non-cardinality finite-limit measure.} Try Jukna's corrected switching/finite-limit criterion with a Hamming or support-sensitive length measure. Falsifier: the requirement that selected positives respect the measure may force the measure back to cardinality on every small set.
\paragraph{B: connected cumulant potential.} Reorganize the factorial-moment inclusion--exclusion series into factorial cumulants/connected correlations and test whether low-order connected terms already collapse. Falsifier: an easy family with exponential Bonferroni cancellation but a trivial decision bit.
\paragraph{C: bounded cancellation from small set circuits.} Seek a theorem that a small union/intersection DAG can only realize bounded-order cancellation. Falsifier: mass production or an easy promise family whose witness-count series has huge cancellation depth despite a tiny decision circuit.

\section{LITERATURE CHECK AND REVISIONS}
Jukna's 2018 correction to the monotone switching lemma explicitly requires selected positive and negative inputs to \emph{respect} the chosen length measures. For a measure $\mu$, a selected input $F$ respects $\mu$ iff for every $T\subseteq F$, any coordinate $x$ with $\mu(T\cup\{x\})=\mu(T)$ already belongs to $F$. This condition is load-bearing and is exactly what Candidate A must satisfy.

Factorial-cumulant generating functions use
\[G(z)=\mathbb E[(1+z)^Z]=\sum_{r\ge0}\frac{\mathbb E[(Z)_r]}{r!}z^r,\qquad \log G(z)=\sum_{r\ge1}\frac{f_r}{r!}z^r.\]
Standard cluster/Ursell expansions reorganize inclusion--exclusion into connected correlations when an appropriate convergence condition holds. Such reorganization can be dramatically shorter than raw Bonferroni cancellation, so Candidate B needs an independent high-order cumulant theorem rather than merely large ordinary moments.

Hiltgen--Paterson mass production and the project's UNIVERSALDAG controls remain decisive warnings for Candidate C: a small DAG can share exponentially many implicit subcomputations. Therefore any cancellation-order measure must be proved gate-stable rather than inferred from an expanded formula.

\section{CONSTRUCTION OR PROOF ATTEMPT}
\subsection{A general respect-closure theorem}
Let $X$ be the clause-coordinate universe and $\mathcal A\subseteq2^X$ a selected positive family. Define its forcing closure
\[\operatorname{cl}_{\mathcal A}(T)=\bigcap\{F\in\mathcal A:T\subseteq F\},\]
whenever the indexing family is nonempty.

\begin{theorem}[Respect implies semantic closure]
Let $\mu:2^X\to\mathbb N$ be a Jukna length measure and suppose every $F\in\mathcal A$ respects $\mu$. If $T$ is contained in at least one member of $\mathcal A$, then
\[\{x:\mu(T\cup\{x\})=\mu(T)\}\subseteq\operatorname{cl}_{\mathcal A}(T).\]
\end{theorem}
\begin{proof}
If $\mu(T\cup\{x\})=\mu(T)$ and $F\in\mathcal A$ contains $T$, respect of $F$ forces $x\in F$. This holds for every such $F$, hence $x$ is in their intersection.
\end{proof}

Thus a non-cardinality measure can have a zero-cost extension only when that coordinate is \emph{forced} by every selected positive continuation of the current partial set. Pseudorandom positive families with trivial forcing closure therefore force cardinality-like measures.

\subsection{Cardinality collapse for the full UNSAT slice}
Let $N=8\binom q3$, $t=3q$, and let $\mathcal U_{q,t}$ be all exact-$t$ UNSAT proper 3-CNFs.

\begin{lemma}[UNSAT extension richness]
If $|T|\le t-8$, $x\notin T$, and $\binom q3>|T|+1$, then there exists $F\in\mathcal U_{q,t}$ with $T\subseteq F$ and $x\notin F$.
\end{lemma}
\begin{proof}
Choose a variable triple whose support is used by neither $x$ nor any clause of $T$. Add all eight signed clauses on that support; they form an UNSAT core. Since $|T|+8\le t$ and $N-1\ge t$, fill arbitrarily to exactly $t$ distinct clauses while excluding $x$.
\end{proof}

\begin{theorem}[Respect-collapse]
If every $F\in\mathcal U_{q,t}$ respects an integer length measure $\mu$ with $\mu(\emptyset)=0$, $\mu(S)\le|S|$, and monotonicity, then
\[\boxed{\mu(S)=|S|\quad\text{for every }|S|\le t-7}\]
for all sufficiently large $q$.
\end{theorem}
\begin{proof}
Induct on $s=|S|$. Assume every $s$-set $T$, $s\le t-8$, has $\mu(T)=s$. For $x\notin T$, the extension-richness lemma supplies an UNSAT exact-$t$ formula $F\supseteq T$ excluding $x$. Respect forbids $\mu(T\cup\{x\})=\mu(T)$, so integrality gives $\mu(T\cup\{x\})\ge s+1$. The upper bound by cardinality gives equality. The base case is $\mu(\emptyset)=0$.
\end{proof}

This substantially narrows Candidate A. Any useful non-cardinality finite-limit measure must restrict to a much smaller positive subfamily having genuine deterministic closure. But rare deterministic closure is precisely the kind of local fingerprint that earlier slice experiments found dangerous. The full positive family offers no non-cardinality gain through essentially its entire relevant size range.

\subsection{Exponential Bonferroni cancellation depth}
Let $F$ be uniform on the exact-$t$ layer and let
\[Z(F)=|\{a\in\{0,1\}^q:a\models F\}|.\]
Write
\[M_r=\mathbb E[(Z)_r],\qquad a_r=M_r/r!.\]
Then
\[\Pr[Z=0]=\sum_{r=0}^{2^q}(-1)^r a_r.\]

\begin{lemma}[Factorial-moment growth recurrence]
For every $r\ge0$,
\[M_{r+1}\ge(M_1-r)M_r.\]
\end{lemma}
\begin{proof}
The functions $z\mapsto(z)_r$ and $z\mapsto z$ are increasing on nonnegative integers. For an independent copy $Z'$,
\[\operatorname{Cov}((Z)_r,Z)=\tfrac12\mathbb E[((Z)_r-(Z')_r)(Z-Z')]\ge0.\]
Hence $\mathbb E[Z(Z)_r]\ge M_1M_r$. Since $Z(Z)_r=(Z)_{r+1}+r(Z)_r$, the claim follows.
\end{proof}

\begin{corollary}[Increasing Bonferroni terms]
If $2r+1\le M_1$, then $a_{r+1}\ge a_r$.
\end{corollary}

\begin{theorem}[Exponential cancellation depth]
Let
\[R_q=\left\lfloor\frac{M_1-1}{2}\right\rfloor.\]
For every truncation order $m\le R_q+1$, the Bonferroni partial sum
\[B_m=\sum_{r=0}^m(-1)^ra_r\]
is trivial in the sense that
\[B_m\ge1\quad(m\text{ even}),\qquad B_m\le1-M_1<0\quad(m\text{ odd}).\]
Moreover
\[M_1=2^q\frac{\binom{7\binom q3}{3q}}{\binom{8\binom q3}{3q}}=\exp((\lambda+o(1))q),\]
where
\[\lambda=\log2+3\log(7/8)=\log(686/512)\approx0.2925530027.\]
Therefore
\[\boxed{R_q=\exp((0.2925530027+o(1))q).}\]
\end{theorem}
\begin{proof}
The recurrence makes $a_0,a_1,\ldots,a_{R_q+1}$ nondecreasing. Pair adjacent terms. Even partial sums are $1+(a_2-a_1)+\cdots\ge1$; odd partial sums are $(1-a_1)+(a_2-a_3)+\cdots\le1-a_1$. The formula for $M_1$ is exact because a fixed assignment satisfies exactly $7\binom q3$ of the $8\binom q3$ proper clauses. Expanding the logarithm of the hypergeometric ratio gives the stated exponent with error $O(1/q)$.
\end{proof}

Thus ordinary inclusion--exclusion does not yield even a nontrivial probability bound until exponentially high order.

\subsection{Connected-cumulant check and a strict second-order effect}
Let $f_r$ be the factorial cumulants defined by
\[\log G(z)=\sum_{r\ge1}f_rz^r/r!,\qquad G(z)=\mathbb E[(1+z)^Z].\]
Then $f_1=M_1$ and $f_2=M_2-M_1^2$.

Consider ordered assignment pairs at relative Hamming distance $\delta=1/3$. Their number has exponential rate $\log2+h(1/3)$, and the fraction of clauses satisfied by both assignments is
\[s_2(1/3)=\frac{6+(2/3)^3}{8}=\frac{85}{108}.\]
Hence
\[\liminf_{q\to\infty}\frac1q\log M_2\ge c_2:=\log2+h(1/3)+3\log(85/108).\]
The exponent gap is strictly positive:
\[c_2-2\lambda>0,\]
because after exponentiating three times this is equivalent to
\[497393604742873088000000000>459912097657205621216589969.\]
Therefore
\[\boxed{f_2=M_2(1-o(1))\text{ on the exponential scale, and }f_2\ge\exp((c_2-o(1))q).}\]
So the witness-count law is not asymptotically Poisson even at connected order two.

Numerical method-of-types optimization for fixed orders $r=1,\ldots,5$ gave candidate exponents
\[0.2925530,\ 0.6118424,\ 0.9537562,\ 1.3141773,\ 1.6899878.\]
Each exceeds every sum of the earlier candidate exponents over a nontrivial partition of $r$. This suggests the concrete conjecture
\[\Lambda_r>\max_{\pi\text{ nontrivial}}\sum_{B\in\pi}\Lambda_{|B|},\]
which would force every fixed-order factorial cumulant to have the same leading exponent as its moment. This higher-order statement is \COMP/OPEN, not promoted.

\subsection{Why cancellation depth alone cannot lower-bound decision circuits}
Bonferroni depth is a property of a chosen witness-event expansion, not of the one-bit decision function. An extreme control is a random variable $Z\in\{0,M\}$ whose zero/nonzero status is supplied by one source bit. Its decision circuit has size one, while its inclusion--exclusion series has factorial moments $M_r=\tfrac12(M)_r$ and requires order $\Theta(M)$ before terms stop increasing.

Equivalently, in a proper-CNF promise one may choose one fixed UNSAT exact-$t$ formula and one SAT exact-$t$ formula differing in a designated core clause; on this two-input promise one coordinate distinguishes them, while the SAT member can have exponentially many satisfying assignments. Thus exponential witness-cancellation depth does not imply decision hardness without a theorem tying the chosen witness expansion to every legal circuit.

A second hostile control is Poisson: if $Z$ were Poisson with mean $\mu$, raw inclusion--exclusion requires order about $\mu$, but
\[\log G(z)=\mu z\]
has only one nonzero factorial cumulant. Hence the cumulant reorganization can collapse enormous Bonferroni cancellation. Our positive $f_2$ theorem rules out this exact Poisson collapse for random 3-SAT, but not a bounded higher-order collapse.

\section{TESTS AND OUTCOME}
The arithmetic audit evaluated the exact mean witness count at $q=8,16,32,64,128$, obtaining approximately $9.478$, $104.022$, $1.1444\times10^4$, $1.3429\times10^8$, and $1.8250\times10^{16}$. The corresponding guaranteed trivial-Bonferroni ranges $(M_1-1)/2$ are approximately $4.24$, $51.5$, $5.72\times10^3$, $6.71\times10^7$, and $9.12\times10^{15}$. The fixed-order cumulant-rate optimizations through $r=5$ were separately checkpointed. These are \COMP confirmations of the symbolic statements.

\paragraph{Outcome.}
Candidate A is \RET for the full exact-density UNSAT positive family: respect forces cardinality through size $t-7$. Candidate B survives only in the sharpened form of a high-order connected-cumulant theorem; ordinary cancellation depth is insufficient. Candidate C is \RET in its generic form: no decision lower bound follows from witness-series cancellation without a gate-stable transfer. The highest-value next targets are (i) prove strict cumulant indecomposability for order growing with $q$, and separately (ii) find a set-circuit invariant whose union/intersection behavior forces such connected order to grow with circuit size.

Root \texttt{main.tex} remains truncated and was not modified.

\begin{thebibliography}{6}
\bibitem{Jukna18} S. Jukna, \emph{Corrected Lower-Bounds Criterion and Monotone Switching Lemma}, online notes to Theorem 9.26, 2018.
\bibitem{Jukna97} S. Jukna, \emph{Finite Limits and Monotone Computations: The Lower Bounds Criterion}, CCC 1997.
\bibitem{Jansen19} S. Jansen, \emph{Cluster Expansions for Gibbs Point Processes}, Advances in Applied Probability 51, 2019.
\bibitem{HP95} A. P. Hiltgen and M. S. Paterson, \emph{PI$_k$ Mass Production and an Optimal Circuit for the Ne\v{c}iporuk Slice}, Computational Complexity 5, 1995.
\bibitem{Wegener85} I. Wegener, \emph{On the Complexity of Slice Functions}, TCS 38, 1985.
\bibitem{HS03} M. T. Hajiaghayi and G. B. Sorkin, \emph{The Satisfiability Threshold of Random 3-SAT Is at Least 3.52}, 2003.
\end{thebibliography}
\end{document}
% FINAL_VERIFICATION_SENTINEL::MAIN::SLICECANCELDEPTH1::2026-10-18_0205::P_VS_NP_OPEN